\section{Document acquisition and classification}
\label{sec:acquisition}

HMRS parses PDFs, scanned images, Word and Excel files through
an asynchronous MinerU service \citep{mineru}. The MCP parse tools return
pollable task identifiers. Markdown and plain text take the direct save
path. For parsed inputs, the ingest agent reviews the recovered text
before selecting a collection by topic. The collection, document and segment
levels are listed in Table~\ref{tab:levels}; Figure~\ref{fig:levels}
illustrates the catalogue-to-passage path after material has been stored.

The ingestion workflow connects parsing, review, optional splitting,
placement, storage and indexing. Depending on the operations performed,
the workflow may produce a parse task identifier, quality report, content
fingerprint, split manifest, storage record and destination collection; these
are not uniform platform receipts. After storage and indexing, a content-based retrieval probe
checks whether the saved material is searchable.

\subsection{Quality gates and ingest lifecycle}

The upload service checks sibling files for exact duplicates using content
hashes. A separate near-duplicate tool returns pairs above a similarity
threshold for review. The ingest agent checks recovered text for
optical character recognition (OCR) noise,
binary residue, missing section bodies, whitespace collapse and language
mismatch. It selects a collection whose scope matches the recovered subject and
derives descriptions and tags from the same body. The owning collection and
its listing can then be checked after storage.

For oversized documents, the optional structural splitter groups recognised
units such as headings, Markdown tables, lists, block quotes and fenced code.
The arXiv campaign used a configured window of 30,000 characters; later configurations used
60,000. The splitter aims to fill 85\% of that window and records
oversized indivisible units where applicable. Structure recognition covers
Markdown units, with formula-heavy pages depending on the quality of the
parsed representation.

The web and MCP entry points share parsing and storage, but differ at
indexing. The web parsed-file route saves Markdown and images; indexing
is separate. The MCP save tool can schedule vector and lexical indexing in
the background. Graph construction is separate in both cases. The reference
workflow treats material as retrieval-ready only after the relevant indexing
operation and a successful content-based probe. Post-index checks compare
expected identifiers with recorded chunks and query for saved content.
Figure~\ref{fig:arch} separates managed files and metadata from their
index and graph projections.
